\section{Introduction}\label{sec:intro}

\begin{contract}
\textbf{Assumes:} nothing. \quad
\textbf{Produces:} the statement of the conjecture and of the main theorem
(\cref{thm:main}), the two elementary observations that fix the meaning of the definitions, and a
map of the proof (\cref{fig:roadmap}) naming its three engines.\\[2pt]
\textbf{Why:} the proof is an interplay of three ingredients that are individually simple and
collectively confusing. Knowing which is which before starting is most of the battle.
\end{contract}

\subsection{The conjecture}

An \emph{oriented graph} is a finite directed graph with no loops, no two arcs with the same tail
and head, and no pair of opposite arcs: if $x\to y$ then not $y\to x$. Its underlying undirected
graph need not be complete and need not be connected.

For a vertex $v$,
\begin{equation}\label{eq:N1N2}
 N_1^+(v)=\{u:\ v\to u\},\qquad
 N_2^+(v)=\bigl\{w\notin\{v\}\cup N_1^+(v):\ \text{there is }u\text{ with }v\to u\to w\bigr\}.
\end{equation}
So $N_1^+(v)$ is the set of out-neighbors and $N_2^+(v)$ is the set of vertices at directed
distance \emph{exactly} two. Each vertex is counted once, however many two-step paths reach it.

\begin{figure}[t]
\centering
\begin{tikzpicture}[
  vertex/.style={circle,draw=black!70,fill=black!6,inner sep=0pt,minimum size=5.6mm,
                 font=\footnotesize},
  first/.style={circle,draw=blue!55!black,fill=blue!12,inner sep=0pt,minimum size=5.6mm,
                font=\footnotesize},
  second/.style={circle,draw=green!45!black,fill=green!14,inner sep=0pt,minimum size=5.6mm,
                 font=\footnotesize},
  far/.style={circle,draw=black!35,fill=white,inner sep=0pt,minimum size=5.6mm,
              font=\footnotesize,text=black!55},
  arc/.style={-{Stealth[length=1.6mm]},draw=black!75},
  lab/.style={font=\scriptsize\itshape,text=black!70}
]
% vertices
\node[vertex] (v) at (0,0) {$v$};
\node[first]  (a) at (1.75,1.55) {$a$};
\node[first]  (b) at (2.05,0)    {$b$};
\node[first]  (c) at (1.75,-1.55){$c$};
\node[second] (p) at (4.05,2.55) {$p$};
\node[second] (q) at (4.35,1.35) {$q$};
\node[second] (r) at (4.45,-0.25){$r$};
\node[second] (s) at (4.15,-1.85){$s$};
\node[far]    (z) at (6.05,2.05) {$z$};
% arcs
\draw[arc] (v) -- (a);
\draw[arc] (v) -- (b);
\draw[arc] (v) -- (c);
\draw[arc] (a) -- (p);
\draw[arc] (a) -- (q);
\draw[arc] (b) -- (r);
\draw[arc] (c) -- (s);
\draw[arc] (c) to[bend left=18] (b);
\draw[arc] (p) -- (z);
% region labels (placed clear of every node and arc)
\node[lab,anchor=west] at (1.35,2.35) {$N_1^+(v)$};
\node[lab,anchor=west] at (4.55,3.15) {$N_2^+(v)$};
\node[lab,anchor=west,text=black!55] at (6.25,1.55) {distance $3$};
\node[lab,anchor=north] at (2.85,-2.55) {$b$ is reachable in one \emph{and} in two steps,
                                        so $b\notin N_2^+(v)$};
\end{tikzpicture}

\caption{The two neighborhoods of \eqref{eq:N1N2}, and the two things that can go wrong if one is
 careless. \emph{Look at} the vertex $b$: it lies in $N_1^+(v)$ and is also reachable in two steps
 ($v\to c\to b$), but it is \emph{not} in $N_2^+(v)$, because distance two must be exact. Look
 also at $z$, reachable only in three steps. \emph{Why it matters:} the subtraction of $N_1^+(v)$
 in \eqref{eq:N1N2} is what makes the conjecture a comparison of two disjoint sets, and the first
 place where a proof can fail by forgetting a case.}
\label{fig:definitions}
\end{figure}

\begin{quote}\itshape
Seymour's second-neighborhood conjecture. Every nonempty finite oriented graph has a vertex $v$
with $\abs{N_1^+(v)}\le\abs{N_2^+(v)}$.
\end{quote}

\begin{theorem}[Main theorem]\label{thm:main}
Every nonempty finite oriented graph has a vertex $v$ such that
$\abs{N_1^+(v)}\le\abs{N_2^+(v)}$.
\end{theorem}

\inwords{In every finite oriented graph, some vertex has at least as many vertices at distance
exactly two as at distance exactly one.}

\begin{proof}[Proof]
Combine \cref{prop:reduction} with \cref{prop:no-concavity}: the first turns a counterexample into
a graph of a certain kind, the second shows that no graph of that kind exists. The two propositions
are proved in \cref{sec:reduction} and \cref{sec:incompatibility}; \cref{fig:roadmap} shows the
dependencies.
\end{proof}

\Cref{thm:main} is unconditional: no connectivity, no minimum degree, no regularity. This is not a
defect of the statement --- a sink (a vertex with no out-neighbor) satisfies the conclusion
trivially, with $0\le 0$ --- but it does mean that the interesting case, and the whole difficulty,
is a graph in which every vertex has both in- and out-neighbors.

\begin{warnbox}
Two natural weakenings are \emph{not} what is proved here, and the reader should not substitute
them. First, replacing $N_2^+$ by $F^2(\{v\})$ without subtracting $N_1^+(v)$ would make the
statement false in general digraphs and trivial in tournaments. Second, the theorem is a
\emph{maximum over vertices} of the ratio $\abs{N_2^+}/\abs{N_1^+}$, not an average: the average can
be smaller than $1$ (take a long directed path).
\end{warnbox}

\subsection{The shape of the proof: three inventions}

The proof is a single argument, but three of its steps are genuinely new and everything else is
either a reduction to them or a consequence of them. They are worth naming at the start, because
each is a different kind of idea.

\begin{ideabox}[title={The three engines}]
\begin{enumerate}[leftmargin=1.3em,itemsep=2pt,topsep=1pt]
\item \textbf{A subset inequality with the right right-hand side}
  (\cref{lem:deficit}). In a minimal counterexample, deleting a well-chosen set of arcs and using
  minimality shows
  \[
    \abs{F^2(S)\setminus F(S)}<\abs{F(S)\setminus S}
    \qquad\text{for every nonempty proper }S .
  \]
  The point is the right-hand side: $\abs{F(S)\setminus S}$, the part of the first image
  \emph{outside} $S$, rather than the whole first image. A transitive-tournament blowup
  (\cref{lem:blowup}) then turns this strict deficit into
  $\abs{U}+\abs{F^2(U)}<2\abs{F(U)}$.

\item \textbf{Pruning at a cost charged to what stays uncovered} (\cref{lem:pruning}). Two families
  of ordered pairs that are \emph{almost} compatible can be made compatible by deleting a few
  members; the number deleted is bounded by the number of collisions plus the number of points of
  a certain set $Z$ that remain uncovered \emph{after} the pruning. Charging the cost to the
  survivors is what makes the last step work.

\item \textbf{An extremal objective that refunds the cost} (\cref{sec:incompatibility}). Maximise
  $\abs{P}+\abs{Q}$ plus the number $u$ of uncovered points over all compatible pairs of families
  containing the diagonal. The subset inequality gives a pair worth more than $M+2u$; pruning it
  costs at most $u+u'$ and leaves $u'$ uncovered, so the net value exceeds $M+u$ --- a
  contradiction.
\end{enumerate}
\end{ideabox}

\begin{figure}[t]
\centering
\begin{tikzpicture}[
  box/.style={draw=black!55,fill=white,rounded corners=2pt,align=center,
              font=\scriptsize,inner sep=3pt,text width=#1},
  box/.default=5.2cm,
  thm/.style={box=8.2cm,draw=black!75,fill=black!6,font=\small},
  ar/.style={-{Stealth[length=1.8mm]},draw=black!70,thick}
]
\node[thm] (T) at (0,0) {Main Theorem: some $v$ has $\abs{N_1^+(v)}\le\abs{N_2^+(v)}$};
\node[box,draw=blue!45!black,fill=blue!6] (P2) at (-4.3,-2.1)
  {\textbf{\cref{prop:reduction}}\\[1pt] a counterexample yields $G$ with
   $\abs{U}+\abs{F^2(U)}<2\abs{F(U)}$ for all $\varnothing\ne U\subsetneq V(G)$,
   every indegree positive};
\node[box,draw=blue!45!black,fill=blue!6] (P4) at (4.3,-1.75)
  {\textbf{\cref{prop:no-concavity}}\\[1pt] no such $G$ exists};
\node[box] (L22) at (-4.3,-4.15)
  {\textbf{\cref{lem:deficit}} (subset deficit)\\[1pt]
   $\abs{F^2(S)\setminus F(S)}<\abs{F(S)\setminus S}$ in a minimal counterexample};
\node[box] (L23) at (-4.3,-5.9)
  {\textbf{\cref{lem:blowup}} (blowup)\\[1pt] turns the deficit into the strict inequality};
\node[box] (L31) at (4.3,-3.5)
  {\textbf{\cref{lem:pruning}} (pruning)\\[1pt] two almost compatible families become compatible};
\node[box] (EX) at (4.3,-5.25)
  {\textbf{extremal choice} (Section 4)\\[1pt] maximise $\abs P+\abs Q+u$ over compatible
   pairs containing the diagonal};
\node[box] (L32) at (4.3,-7.1)
  {\textbf{\cref{lem:kernel}} + K\"onig\\[1pt] ranks bound a matching, a matching bounds a cover};
\draw[ar] (P2.north) -- (T.south -| P2.north);
\draw[ar] (P4.north) -- (T.south -| P4.north);
\draw[ar] (L22.north) -- (P2.south);
\draw[ar] (L23.north) -- (L22.south);
\draw[ar] (L31.north) -- (P4.south);
\draw[ar] (EX.north) -- (L31.south);
\draw[ar] (L32.north) -- (EX.south);
\end{tikzpicture}

\caption{The architecture of the proof, read from the bottom up. The main theorem rests on two
 propositions: one converts a counterexample into a graph satisfying a strict subset inequality,
 the other forbids such a graph. Each proposition in turn rests on one substantial lemma.
 \emph{Look at} the two lower boxes: the extremal choice sits between the pruning lemma and the
 contradiction, and is where the two halves of the argument meet. \emph{Why it matters:} \cref{sec:weighted,sec:cycles} are not in this diagram: they consume the
 theorem and produce nothing that is used earlier, so they can be skipped entirely on a first
 reading.}
\label{fig:roadmap}
\end{figure}

\paragraph{How the reduction goes, in words.}
Suppose the theorem fails and take a counterexample $D$ of minimum order, and among those of
minimum number of arcs. Deleting a cleverly chosen family of arcs produces a smaller graph, which
by minimality is \emph{not} a counterexample: some vertex $v$ of it satisfies the desired
inequality. Comparing the neighborhoods of $v$ before and after the deletion shows that the
deleted arcs ``bought'' more second out-neighbors for $v$ than first out-neighbors. Counting this
gain and loss over a chain of deletions, for an arbitrary subset $S$ of vertices, gives the subset
deficit of \cref{lem:deficit}. The blowup then amplifies an integer deficit (at least $1$) into a
margin large enough to survive the final rounding.

\paragraph{How the contradiction goes, in words.}
Represent a \emph{family of ordered pairs} $P\subseteq X\times X$ as a relation: it has a
\emph{column} $P_j=\{p:(p,j)\in P\}$ for each second coordinate and a \emph{row} for each first
coordinate. Advancing the first coordinate along an arc maps $P$ to $\Phi_1P$; advancing the second
coordinate maps a second family $Q$ to $\Phi_2Q$. The families are \emph{compatible} when these two
images are disjoint. Summing the subset inequality over the columns of $P$ and the rows of $Q$ turns
compatibility into a strict inequality between $\abs{P}+\abs{\Phi_1^2P}$ and $2\abs{\Phi_1P}$; taking
complements of the \emph{second} images then produces two families that are worth more than
$M+2u$ but may conflict. The pruning lemma repairs the conflicts at a cost that the uncovered
points pay for. The maximality of $(P,Q)$ is contradicted.

\begin{checkpoint}
After this section you should be able to say, without symbols: \emph{a minimal counterexample
satisfies a strict subset inequality; two families of pairs get pruned to compatibility at a
refundable cost; an extremal pair cannot survive being improved.} Everything that follows is the
proof of one of those three sentences.
\end{checkpoint}

\subsection{History and context}\label{sec:history}

The conjecture was recorded by Dean and Latka~\cite{DeanLatka1995}. The case of
\emph{tournaments} --- oriented graphs with an arc between every two distinct vertices, the case
also known as Dean's conjecture --- was proved by Fisher~\cite{Fisher1996}, using a probability
distribution obtained from Farkas' lemma. Havet and Thomass\'e gave a combinatorial proof using
local median orders~\cite{HavetThomasse2000}. Fidler and Yuster~\cite{FidlerYuster2007} extended
the tournament theorem to tournaments missing a matching, a star, or a clique; Ghazal corrected the
star case and treated generalized stars~\cite{Ghazal2012}. Later extensions cover a matching
together with a star~\cite{DaraFrancisJacobNarayanan2022} and two
stars~\cite{DaamouchAlMninyGhazal2025}.

A second line imposes degree conditions or weakens the ratio. Kaneko and Locke settled the case of
minimum out-degree at most six~\cite{KanekoLocke2001}, and a recent computer-assisted preprint
treats minimum out-degree seven~\cite{SadhukhanSandeepSen2026}. Chen, Shen and
Yuster~\cite{ChenShenYuster2003} proved a universal ratio second-to-first of about $0.657298$;
Huang and Peng~\cite{HuangPeng2024} raise this to about $0.715538$ and report a computational bound
of $0.74530$. Botler, Moura and Naia~\cite{BotlerMouraNaia2023}, followed by Espuny D\'iaz,
Gir\~ao, Granet and Kronenberg~\cite{EspunyDiazEtAl2025}, treated every orientation of a binomial
random graph; the latter covers every fixed edge probability below $1/2$ with probability tending
to one. Brukhman's recent dense-case preprint~\cite{Brukhman2026} concerns order $n\le 2\delta+2$,
where $\delta$ is the minimum out-degree.

The present proof builds on the minimal-counterexample and graph-product methods of Brantner,
Brockman, Kay and Snively~\cite{BrantnerBrockmanKaySnively2009} and on Seacrest's subset
perspective~\cite{Seacrest2019}. It adapts the edge-deletion idea in the proof of Seacrest's
Lemma~4 --- specializing his parameter to $1$ --- to obtain the stronger subset inequality needed
here, and gives that argument in full. The main additional ingredient is the pruning bound for
arbitrary finite binary relations (\cref{lem:pruning}), whose deletion cost depends on the points
left uncovered after pruning.

\subsection{Organization of the paper}\label{sec:organisation}

\begin{description}[leftmargin=1.4em,itemsep=2pt]
\item[\Cref{sec:reduction}] From a counterexample to a strict subset inequality
  (\cref{prop:reduction}), via the subset deficit (\cref{lem:deficit}) and the blowup
  (\cref{lem:blowup}). This is the \emph{first engine} and the easiest of the three to read.
\item[\Cref{sec:pruning}] The pruning lemma (\cref{lem:pruning}) for arbitrary finite binary
  relations, and the rank lemma (\cref{lem:kernel}) that makes it work. This is the \emph{second
  engine}; it is independent of graph theory and can be read on its own.
\item[\Cref{sec:incompatibility}] The extremal contradiction (\cref{prop:no-concavity}). This is
  the \emph{third engine} and the place where the first two meet.
\item[\Cref{sec:weighted,sec:cycles}] Consequences: the vertex- and arc-weighted forms via
  Seacrest's equivalences, and consequences for directed triangles, minimum degree and girth.
  Neither is used in the proof of \cref{thm:main}.
\item[\appref{app:toolkit}] The five facts used from outside graph theory, with proofs: K\H{o}nig's
  theorem, rank--nullity, Sylvester's inequality, the minor--matching bound, and the Farkas
  alternative. Nothing in the main text depends on anything else.
\item[\appref{app:alternatives}] What was tried in search of a shorter proof: two simplifications that
  were adopted (no Edmonds, no Seacrest for the vertex-weighted forms) and four attempts that failed,
  with the reason.
\item[Appendices~\ref{app:depgraph}, \ref{app:dictionary} and \ref{app:checks}] A dependency graph of the numbered statements,
  a dictionary of notation against the original preprint, and a record of how each step was
  checked (re-derivation, exhaustive computation on small cases, and a Lean formalisation).
\end{description}
